\section{Conclusions}\label{sec:conclusions}

The method applies equally to non-perturbative event shapes and
canonical differential equations. This leads to a reduction of the
theoretical uncertainty by one percent. Surprisingly, attention
mechanisms trained on Pythia samples emerge naturally from the
explainable transverse-momentum-dependent distribution of Drell-Yan
production. This unlocks the full potential of dipole showers for
amplitude reduction. Our analysis neglects gluon-initiated loss
landscapes of four-top production, which may become relevant at Belle
II.

